\chapter{Die Mogiljanskaja-Familien}

\section{Die beiden Grundfamilien}

Für die Konstruktion zweier unendlicher Halbgruppen werden zwei
disjunkte getaggte Schichten benötigt. Dieses Kapitel baut ihre familien-
und mengentheoretischen Eigenschaften auf. Die Operationen auf den
Trägermengen werden im folgenden Kapitel definiert. Die allgemeine Theorie
indizierter Familien wird aus Band~21 vorausgesetzt.
Die Zeichen \(0\) und \(1\) in den folgenden geordneten Paaren sind dabei
ausdrücklich die natürlichen Peano-Zahlen und dienen lediglich als Tags.

\MogStatementTarget{MogiljanskajaLayerLDef}
\FormulaDefDeltaK[Die einfach indizierte Grundschicht]{%
  L\coloneqq\{0\}\times\mathbb N%
}{MogiljanskajaLayerLDef}{%
  \DeltaRow{Mengen}{L}
}

\MogStatementTarget{MogiljanskajaLayerDDef}
\FormulaDefDeltaK[Die zweifach indizierte Grundschicht]{%
  D\coloneqq\{1\}\times(\mathbb N\times\mathbb N)%
}{MogiljanskajaLayerDDef}{%
  \DeltaRow{Mengen}{D}
}

\MogStatementTarget{MogiljanskajaAElementsDef}
\FormulaDefDeltaK[Die einfach indizierte Grundfamilie]{%
  \begin{aligned}[t]
    &a\colon\mathbb N\to L,\\[-2pt]
    &\forall i\in\mathbb N\,\bigl(a_i\coloneqq(0,i)\bigr)
  \end{aligned}%
}{MogiljanskajaAElementsDef}{%
  \DeltaRow{Mengen}{L\dsep a_i\dsep i}
  \DeltaRow{Funktionen}{a}
}

\begin{tabproofsplitwide}
  \proofpartwide{Typisierung der Grundfamilie}
    \proofstepwidestar[1]{%
      i\in\mathbb N}{%
      \rA}
    \proofstepwidestar[1]{%
      (0,i)\in\{0\}\times\mathbb N}{%
      \FormulaRefAuto{a\in A\dsep b\in B\vdash(a,b)\in A\times B}[thm]{\FormulaRefAuto{a\in\{a\}}[thm],1}}
    \proofstepwidestar[]{%
      L=\{0\}\times\mathbb N}{%
      \FormulaRefAuto{MogiljanskajaLayerLDef}[def]}
    \proofstepwidestar[1]{%
      (0,i)\in L}{%
      \rIE{3,2}}
    \proofstepwidestar[]{%
      \forall i\in\mathbb N\,((0,i)\in L)}{%
      \rUI{\rRI{1,4}}}
  \closeproofpartwide
\end{tabproofsplitwide}

\MogStatementTarget{MogiljanskajaBElementsDef}
\FormulaDefDeltaK[Die zweifach indizierte Grundfamilie]{%
  \begin{aligned}[t]
    &b\colon\mathbb N\times\mathbb N\to D,\\[-2pt]
    &\forall i,j\in\mathbb N\,\bigl(b_{ij}\coloneqq(1,(i,j))\bigr)
  \end{aligned}%
}{MogiljanskajaBElementsDef}{%
  \DeltaRow{Mengen}{D\dsep b_{ij}\dsep i\dsep j}
  \DeltaRow{Funktionen}{b}
}

Hier steht \(b_{ij}\) für \(b((i,j))\). Die Familie ist also eine einzige
Funktion auf \(\mathbb N\times\mathbb N\), keine zweistufige Sammlung
voneinander unabhängiger Funktionen.

\MogStatementTarget{MogiljanskajaAIndexUnique}
\FormulaThmDeltaK[Eindeutigkeit der (a)-Indizes]{%
  i,j\in\mathbb N\dsep a_i=a_j\vdash i=j%
}{MogiljanskajaAIndexUnique}{%
  \DeltaRow{Mengen}{a_i\dsep a_j\dsep i\dsep j}
  \DeltaRow{Funktionen}{a}
}
\begin{tabproofsplitwide}
  \proofpartwide{Gleichheit der zweiten Koordinate}
    \proofstepwidestar[1]{%
      i,j\in\mathbb N}{%
      \rA}
    \proofstepwidestar[2]{%
      a_i=a_j}{%
      \rA}
    \proofstepwidestar[1,2]{%
      (0,i)=(0,j)}{%
      \FormulaRefAuto{MogiljanskajaAElementsDef}[def]{1,2}}
    \proofstepwidestar[1,2]{%
      i=j}{%
      \FormulaRefAuto{TaggedPairRightInjective}[thm]{3}}
  \closeproofpartwide
\end{tabproofsplitwide}

\MogStatementTarget{MogiljanskajaBIndexUnique}
\FormulaThmDeltaK[Eindeutigkeit der (b)-Indizes]{%
  i,j,m,n\in\mathbb N\dsep b_{ij}=b_{mn}
  \vdash i=m\land j=n%
}{MogiljanskajaBIndexUnique}{%
  \DeltaRow{Mengen}{b_{ij}\dsep b_{mn}\dsep i\dsep j\dsep m\dsep n}
  \DeltaRow{Funktionen}{b}
}
\begin{tabproofsplitwide}
  \proofpartwide{Gleichheit der beiden inneren Koordinaten}
    \proofstepwidestar[1]{%
      i,j,m,n\in\mathbb N}{%
      \rA}
    \proofstepwidestar[2]{%
      b_{ij}=b_{mn}}{%
      \rA}
    \proofstepwidestar[1,2]{%
      (1,(i,j))=(1,(m,n))}{%
      \FormulaRefAuto{MogiljanskajaBElementsDef}[def]{1,2}}
    \proofstepwidestar[1,2]{%
      i=m\land j=n}{%
      \FormulaRefAuto{TaggedNestedPairInjective}[thm]{3}}
  \closeproofpartwide
\end{tabproofsplitwide}

\MogStatementTarget{MogiljanskajaLayerLMembership}
\FormulaThmDeltaK[Elementkriterium der (a)-Schicht]{%
  s\in L\leftrightarrow\exists i\in\mathbb N\,(s=a_i)%
}{MogiljanskajaLayerLMembership}{%
  \DeltaRow{Mengen}{L\dsep s\dsep a_i\dsep i}
  \DeltaRow{Funktionen}{a}
}
\begin{tabproofsplitwide}
  \proofpartwide{Darstellung der getaggten Schicht}
    \proofstepwidestar[]{%
      L=\{0\}\times\mathbb N}{%
      \FormulaRefAuto{MogiljanskajaLayerLDef}[def]}
    \proofstepwidestar[]{%
      s\in L\leftrightarrow\exists i\in\mathbb N\,(s=(0,i))}{%
      \FormulaRefAuto{TaggedProductLayerRepresentation}[thm],\ \rIE{1}}
    \proofstepwidestar[]{%
      \forall i\in\mathbb N\,(a_i=(0,i))}{%
      \FormulaRefAuto{MogiljanskajaAElementsDef}[def]}
    \proofstepwidestar[]{%
      s\in L\leftrightarrow\exists i\in\mathbb N\,(s=a_i)}{%
      2,\ \rUE{3}}
  \closeproofpartwide
\end{tabproofsplitwide}

\MogStatementTarget{MogiljanskajaLayerDMembership}
\FormulaThmDeltaK[Elementkriterium der (b)-Schicht]{%
  s\in D\leftrightarrow
    \exists i,j\in\mathbb N\,(s=b_{ij})%
}{MogiljanskajaLayerDMembership}{%
  \DeltaRow{Mengen}{D\dsep s\dsep b_{ij}\dsep i\dsep j}
  \DeltaRow{Funktionen}{b}
}
\begin{tabproofsplitwide}
  \proofpartwide{Darstellung der doppelt indizierten Schicht}
    \proofstepwidestar[]{%
      D=\{1\}\times(\mathbb N\times\mathbb N)}{%
      \FormulaRefAuto{MogiljanskajaLayerDDef}[def]}
    \proofstepwidestar[]{%
      s\in D\leftrightarrow\exists i,j\in\mathbb N\,(s=(1,(i,j)))}{%
      \FormulaRefAuto{TaggedNestedProductLayerRepresentation}[thm],\ \rIE{1}}
    \proofstepwidestar[]{%
      \forall i,j\in\mathbb N\,(b_{ij}=(1,(i,j)))}{%
      \FormulaRefAuto{MogiljanskajaBElementsDef}[def]}
    \proofstepwidestar[]{%
      s\in D\leftrightarrow\exists i,j\in\mathbb N\,(s=b_{ij})}{%
      2,\ \rUE{\rUE{3}}}
  \closeproofpartwide
\end{tabproofsplitwide}

\MogStatementTarget{MogiljanskajaLayerLImage}
\FormulaThmDeltaK[Bildmenge der (a)-Familie]{%
  a[\mathbb N]=L%
}{MogiljanskajaLayerLImage}{%
  \DeltaRow{Mengen}{L\dsep a[\mathbb N]}
  \DeltaRow{Funktionen}{a}
}
\begin{tabproofsplitwide}
  \proofpartwide{Mengenextensionalität}
    \proofstepwidestar[]{%
      a\colon\mathbb N\to L}{%
      \FormulaRefAuto{MogiljanskajaAElementsDef}[def]}
    \proofstepwidestar[]{%
      s\in a[\mathbb N]\leftrightarrow\exists i\in\mathbb N\,(s=a_i)}{%
      \FormulaRefAuto{IndexedFamilyImageMembership}[thm]{1},\ \FormulaRefAuto{IndexedFamilyImageDef}[def]}
    \proofstepwidestar[]{%
      s\in L\leftrightarrow\exists i\in\mathbb N\,(s=a_i)}{%
      \FormulaRefAuto{MogiljanskajaLayerLMembership}[thm]}
    \proofstepwidestar[]{%
      a[\mathbb N]=L}{%
      \FormulaRefAuto{\forall x\,(x\in A\leftrightarrow x\in B)\vdash A=B}[ax]{\rUI{\FormulaRefAuto{P\leftrightarrow Q,R\leftrightarrow Q\vdash P\leftrightarrow R}[thm]{2,3}}}}
  \closeproofpartwide
\end{tabproofsplitwide}

\MogStatementTarget{MogiljanskajaLayerDImage}
\FormulaThmDeltaK[Bildmenge der (b)-Familie]{%
  b[\mathbb N\times\mathbb N]=D%
}{MogiljanskajaLayerDImage}{%
  \DeltaRow{Mengen}{D\dsep b[\mathbb N\times\mathbb N]}
  \DeltaRow{Funktionen}{b}
}
\begin{tabproofsplitwide}
  \proofpartwide{Mengenextensionalität}
    \proofstepwidestar[]{%
      b\colon\mathbb N\times\mathbb N\to D}{%
      \FormulaRefAuto{MogiljanskajaBElementsDef}[def]}
    \proofstepwidestar[]{%
      s\in b[\mathbb N\times\mathbb N]\leftrightarrow\exists i,j\in\mathbb N\,(s=b_{ij})}{%
      \FormulaRefAuto{IndexedFamilyImageMembership}[thm]{1},\ \FormulaRefAuto{IndexedFamilyImageDef}[def],\ \FormulaRefAuto{x\in A\times B\eqvdash\exists a\in A\,\exists b\in B\,x=(a,b)}[thm]}
    \proofstepwidestar[]{%
      s\in D\leftrightarrow\exists i,j\in\mathbb N\,(s=b_{ij})}{%
      \FormulaRefAuto{MogiljanskajaLayerDMembership}[thm]}
    \proofstepwidestar[]{%
      b[\mathbb N\times\mathbb N]=D}{%
      \FormulaRefAuto{\forall x\,(x\in A\leftrightarrow x\in B)\vdash A=B}[ax]{\rUI{\FormulaRefAuto{P\leftrightarrow Q,R\leftrightarrow Q\vdash P\leftrightarrow R}[thm]{2,3}}}}
  \closeproofpartwide
\end{tabproofsplitwide}

\MogStatementTarget{MogiljanskajaLayersDisjoint}
\FormulaThmDeltaK[Die beiden Grundschichten sind disjunkt]{%
  L\cap D=\varnothing%
}{MogiljanskajaLayersDisjoint}{%
  \DeltaRow{Mengen}{L\dsep D}
}
\begin{tabproofsplitwide}
  \proofpartwide{Verschiedene Tags}
    \proofstepwidestar[]{%
      0\neq1}{%
      \FormulaRefAuto{PeanoZeroNeOne}[thm]}
    \proofstepwidestar[]{%
      (\{0\}\times\mathbb N)\cap(\{1\}\times(\mathbb N\times\mathbb N))=\varnothing}{%
      \FormulaRefAuto{TaggedProductLayersDisjoint}[thm]{1}}
    \proofstepwidestar[]{%
      L\cap D=\varnothing}{%
      \FormulaRefAuto{MogiljanskajaLayerLDef}[def],\ \FormulaRefAuto{MogiljanskajaLayerDDef}[def],\ \rIE{2}}
  \closeproofpartwide
\end{tabproofsplitwide}

\Needspace{20\baselineskip}
\section{Marker und Reserveschichten}

\MogStatementTarget{MogiljanskajaDPrimeDef}
\FormulaDefDeltaK[Marker und unendliche Reserveschicht]{%
  \begin{aligned}[t]
    c_0&\coloneqq b_{00},&
    c_1&\coloneqq b_{11},&
    c_2&\coloneqq b_{10},\\[-2pt]
    U&\coloneqq\bigl\{u\in D\bigm|
      \exists n,j\in\mathbb N\,
        u=b_{(n+1)+1,j}\bigr\},\\[-2pt]
    V&\coloneqq U\cup\{c_2\},&
    P&\coloneqq\{c_0,c_1\},&
    D'&\coloneqq P\cup V
  \end{aligned}%
}{MogiljanskajaDPrimeDef}{%
  \DeltaRow{Mengen}{D\dsep D'\dsep U\dsep V\dsep P\dsep
    c_0\dsep c_1\dsep c_2\dsep b_{ij}\dsep n\dsep j\dsep u}
  \DeltaRow{Funktionen}{b}
}

Somit besteht \(U\) genau aus den Elementen der Zeilen mit erstem Index
mindestens \(2\). Die Menge \(V\) adjungiert den Marker \(c_2=b_{10}\),
und \(D'\) adjungiert zusätzlich \(c_0=b_{00}\) und \(c_1=b_{11}\).

\MogStatementTarget{MogiljanskajaReserveMembership}
\FormulaThmDeltaK[Elementkriterium der Reserveschicht]{%
  u\in U\eqvdash
  \exists n,j\in\mathbb N\,
    (u=b_{(n+1)+1,j})%
}{MogiljanskajaReserveMembership}{%
  \DeltaRow{Mengen}{D\dsep U\dsep u\dsep b_{ij}\dsep n\dsep j}
  \DeltaRow{Funktionen}{b}
}
\begin{tabproofsplitwide}
  \proofpartwide{Die zusätzliche Trägerbedingung ist erfüllt}
    \proofstepwidestar[1]{%
      \exists n,j\in\mathbb N\,(u=b_{(n+1)+1,j})}{%
      \rA}
    \proofstepwidestar[2]{%
      n\in\mathbb N\land(j\in\mathbb N\land u=b_{(n+1)+1,j})}{%
      \rA}
    \proofstepwidestar[2]{n+1\in\mathbb N}{%
      \FormulaRefAuto{PeanoAddOneClosure}[thm]{\rAEa{2}}}
    \proofstepwidestar[2]{%
      (n+1)+1\in\mathbb N}{%
      \FormulaRefAuto{PeanoAddOneClosure}[thm]{3}}
    \proofstepwidestar[2]{%
      b_{(n+1)+1,j}\in D}{%
      \FormulaRefAuto{MogiljanskajaBElementsDef}[def]{4,\rAEa{\rAEb{2}}},\ \FormulaRefAuto{IndexedFamilyCoordinateTyping}[thm]}
    \proofstepwidestar[2]{%
      u\in D}{%
      \rIE{\rAEb{\rAEb{2}},5}}
    \proofstepwidestar[2]{%
      u\in U}{%
      \FormulaRefAuto{MogiljanskajaDPrimeDef}[def]{6,\rEI{\rAI{\rAEa{2},\rEI{\rAEb{2}}}}}}
    \proofstepwidestar[1]{%
      u\in U}{%
      \rEE{1,2,7}}
    \proofstepwidestar[9]{%
      u\in U}{%
      \rA}
    \proofstepwidestar[9]{%
      \exists n,j\in\mathbb N\,(u=b_{(n+1)+1,j})}{%
      \FormulaRefAuto{MogiljanskajaDPrimeDef}[def]{9}}
    \proofstepwidestar[]{%
      u\in U\leftrightarrow\exists n,j\in\mathbb N\,(u=b_{(n+1)+1,j})}{%
      \rLRI{\rRI{9,10},\rRI{1,8}}}
  \closeproofpartwide
\end{tabproofsplitwide}

\MogStatementTarget{MogiljanskajaReserveMarkerFacts}
\hypertarget{mog.statement.B21ReserveAvoidsFirstTwoRows}{}
\hypertarget{mog.statement.B21ReserveMarkerDistinct}{}
\FormulaThmDeltaKR[Die Marker und das ausgesparte Rechteckelement]{%
  \begin{aligned}[t]
    &\text{(i)}\quad c_0,c_1,c_2\in D,\\[-2pt]
    &\text{(ii)}\quad c_0,c_1,c_2\notin U,\\[-2pt]
    &\text{(iii)}\quad c_0,c_1,c_2\text{ paarweise verschieden},\\[-2pt]
    &\text{(iv)}\quad D'\subseteq D,\\[-2pt]
    &\text{(v)}\quad b_{01}\notin D'.
  \end{aligned}%
}{%
  \begin{aligned}[t]
    &c_0,c_1,c_2\in D\dsep
      c_0,c_1,c_2\notin U\dsep
      c_0,c_1,c_2\text{ paarweise verschieden},\\[-2pt]
    &D'\subseteq D\dsep b_{01}\notin D'
  \end{aligned}%
}{MogiljanskajaReserveMarkerFacts}{%
  \DeltaRow{Mengen}{D\dsep D'\dsep U\dsep V\dsep P\dsep
    c_0\dsep c_1\dsep c_2\dsep b_{01}\dsep n\dsep j}
  \DeltaRow{Funktionen}{b}
}
\begin{tabproofsplitwide}
\proofpartwideindR[Die ersten beiden Zeilen liegen außerhalb von \(U\)]{%
    i\in\{0,1\}\dsep j\in\mathbb N\vdash b_{ij}\notin U
  }{i\in\{0,1\}\dsep j\in\mathbb N\vdash b_{ij}\notin U}[B21ReserveAvoidsFirstTwoRows]
    \proofstepwidestar[1]{%
      n\in\mathbb N}{%
      \rA}
    \proofstepwidestar[1]{%
      n+1\in\mathbb N}{%
      \FormulaRefAuto{PeanoAddOneClosure}[thm]{1}}
    \proofstepwidestar[1]{%
      (n+1)+1\neq0}{%
      \FormulaRefAuto{PeanoAddOneNonzero}[thm]{2}}
    \proofstepwidestar[4]{%
      (n+1)+1=1}{%
      \rA}
    \proofstepwidestar[1,4]{0 \in \mathbb N}{%
      \FormulaRefAuto{PeanoZero}[ax]}
    \proofstepwidestar[1,4]{1 \in \mathbb N}{%
      \FormulaRefAuto{PeanoOneInN}[thm]}
    \proofstepwidestar[1,4]{0 + 1 = 1}{%
      \FormulaRefAuto{PeanoAddLeftZero}[thm]{6}}
    \proofstepwidestar[1,4]{%
      n+1=0}{%
      \FormulaRefAuto{PeanoAddOneInjective}[thm]{2,5,%
        \rIE{4,7}}}
    \proofstepwidestar[1,4]{n + 1 \neq 0}{%
      \FormulaRefAuto{PeanoAddOneNonzero}[thm]{1}}
    \proofstepwidestar[1,4]{%
      \bot}{%
      \rBI{9,8}}
    \proofstepwidestar[1]{%
      (n+1)+1\neq1}{%
      \rCI{4,10}}
    \proofstepwidestar[]{%
      \forall n\in\mathbb N\,((n+1)+1\neq0\land(n+1)+1\neq1)}{%
      \rUI{\rRI{1,\rAI{3,11}}}}
    \proofstepwidestar[13]{%
      i\in\{0,1\}\land j\in\mathbb N}{%
      \rA}
    \proofstepwidestar[13]{%
      i\in\mathbb N\land(i=0\lor i=1)}{%
      \FormulaRefAuto{x\in\{a,b\}\eqvdash(x=a\lor x=b)}[thm],\ 13,\ \FormulaRefAuto{PeanoZero}[ax],\ \FormulaRefAuto{PeanoOneInN}[thm]}
    \proofstepwidestar[15]{%
      b_{ij}\in U}{%
      \rA}
    \proofstepwidestar[15]{%
      \exists n,k\in\mathbb N\,(b_{ij}=b_{(n+1)+1,k})}{%
      \FormulaRefAuto{MogiljanskajaReserveMembership}[thm]{15}}
    \proofstepwidestar[17]{%
      n\in\mathbb N\land(k\in\mathbb N\land b_{ij}=b_{(n+1)+1,k})}{%
      \rA}
    \proofstepwidestar[17]{n+1\in\mathbb N}{%
      \FormulaRefAuto{PeanoAddOneClosure}[thm]{\rAEa{17}}}
    \proofstepwidestar[17]{(n+1)+1\in\mathbb N}{%
      \FormulaRefAuto{PeanoAddOneClosure}[thm]{18}}
    \proofstepwidestar[13,17]{i=(n+1)+1\land j=k}{%
      \FormulaRefAuto{MogiljanskajaBIndexUnique}[thm]{14,13,17,19}}
    \proofstepwidestar[13,17]{%
      i=(n+1)+1}{\rAEa{20}
}
    \proofstepwidestar[17]{%
      (n+1)+1\neq0\land(n+1)+1\neq1}{%
      \rRE{\rUE{12},\rAEa{17}}}
    \proofstepwidestar[23]{%
      i=0}{%
      \rA}
    \proofstepwidestar[13,17,23]{%
      (n+1)+1=0}{%
      \rIE{21,23}}
    \proofstepwidestar[13,17,23]{%
      \bot}{%
      \rBI{\rAEa{22},24}}
    \proofstepwidestar[26]{%
      i=1}{%
      \rA}
    \proofstepwidestar[13,17,26]{%
      (n+1)+1=1}{%
      \rIE{21,26}}
    \proofstepwidestar[13,17,26]{%
      \bot}{%
      \rBI{\rAEb{22},27}}
    \proofstepwidestar[13,17]{%
      \bot}{%
      \rOE{\rAEb{14},23,25,26,28}}
    \proofstepwidestar[13,15]{%
      \bot}{%
      \rEE{16,17,29}}
    \proofstepwidestar[13]{%
      b_{ij}\notin U}{%
      \rCI{15,30}}
  \closeproofpartwideindR

  \proofpartwide{Trägereigenschaft}
\proofstepwidestar[]{0\in\mathbb N}{%
      \FormulaRefAuto{PeanoZero}[ax]}
\proofstepwidestar[]{1\in\mathbb N}{%
      \FormulaRefAuto{PeanoOneInN}[thm]}
\proofstepwidestar[]{%
      c_0,c_1,c_2\in D}{%
      \FormulaRefAuto{MogiljanskajaDPrimeDef}[def],\ \FormulaRefAuto{MogiljanskajaBElementsDef}[def]{1,2}}

\closeproofpartwide

\proofpartwide{Erster und zweiter Marker sind verschieden}
  \setcounter{proofstepnr}{3}% Fortlaufende Schritte dieses Beweises.
\proofstepwidestar[4]{%
      c_0=c_1}{%
      \rA}
    \proofstepwidestar[]{0\in\mathbb N}{%
      \FormulaRefAuto{PeanoZero}[ax]}
    \proofstepwidestar[]{1\in\mathbb N}{%
      \FormulaRefAuto{PeanoOneInN}[thm]}
    \proofstepwidestar[4]{b_{00}=b_{11}}{%
      \FormulaRefAuto{MogiljanskajaDPrimeDef}[def]{4}}
    \proofstepwidestar[4]{0=1\land0=1}{%
      \FormulaRefAuto{MogiljanskajaBIndexUnique}[thm]{5,6,7}}
    \proofstepwidestar[4]{%
      0=1}{\rAEa{8}
}
    \proofstepwidestar[4]{0 \neq 1}{%
      \FormulaRefAuto{PeanoZeroNeOne}[thm]}
    \proofstepwidestar[4]{%
      \bot}{%
      \rBI{10,9}}
    \proofstepwidestar[]{%
      c_0\neq c_1}{%
      \rCI{4,11}}

\closeproofpartwide

\proofpartwide{Erster und dritter Marker sind verschieden}
  \setcounter{proofstepnr}{12}% Fortlaufende Schritte dieses Beweises.
\proofstepwidestar[13]{%
      c_0=c_2}{%
      \rA}
    \proofstepwidestar[]{0\in\mathbb N}{%
      \FormulaRefAuto{PeanoZero}[ax]}
    \proofstepwidestar[]{1\in\mathbb N}{%
      \FormulaRefAuto{PeanoOneInN}[thm]}
    \proofstepwidestar[13]{b_{00}=b_{10}}{%
      \FormulaRefAuto{MogiljanskajaDPrimeDef}[def]{13}}
    \proofstepwidestar[13]{0=1\land0=0}{%
      \FormulaRefAuto{MogiljanskajaBIndexUnique}[thm]{14,15,16}}
    \proofstepwidestar[13]{%
      0=1}{\rAEa{17}
}
    \proofstepwidestar[13]{0 \neq 1}{%
      \FormulaRefAuto{PeanoZeroNeOne}[thm]}
    \proofstepwidestar[13]{%
      \bot}{%
      \rBI{19,18}}
    \proofstepwidestar[]{%
      c_0\neq c_2}{%
      \rCI{13,20}}

\closeproofpartwide

\proofpartwide{Zweiter und dritter Marker sind verschieden}
  \setcounter{proofstepnr}{21}% Fortlaufende Schritte dieses Beweises.
\proofstepwidestar[22]{%
      c_1=c_2}{%
      \rA}
    \proofstepwidestar[]{0\in\mathbb N}{%
      \FormulaRefAuto{PeanoZero}[ax]}
    \proofstepwidestar[]{1\in\mathbb N}{%
      \FormulaRefAuto{PeanoOneInN}[thm]}
    \proofstepwidestar[22]{b_{11}=b_{10}}{%
      \FormulaRefAuto{MogiljanskajaDPrimeDef}[def]{22}}
    \proofstepwidestar[22]{1=1\land1=0}{%
      \FormulaRefAuto{MogiljanskajaBIndexUnique}[thm]{23,24,25}}
    \proofstepwidestar[22]{%
      1=0}{\rAEb{26}
}
    \proofstepwidestar[22]{1 \neq 0}{%
      \FormulaRefAuto{PeanoOneNeZero}[thm]}
    \proofstepwidestar[22]{%
      \bot}{%
      \rBI{28,27}}
    \proofstepwidestar[]{%
      c_1\neq c_2}{%
      \rCI{22,29}}

\closeproofpartwide

\proofpartwideindR[Typisierung und Verschiedenheit der Marker]{%
  \begin{aligned}[t]
    &\text{(i)}\quad c_0,c_1,c_2\in D,\\[-2pt]
    &\text{(ii)}\quad c_0\neq c_1\land(c_0\neq c_2\land c_1\neq c_2).
  \end{aligned}%
}{c_0,c_1,c_2\in D\dsep c_0\neq c_1\land(c_0\neq c_2\land c_1\neq c_2)}[B21ReserveMarkerDistinct]
  \setcounter{proofstepnr}{30}% Fortlaufende Schritte dieses Beweises.
\proofstepwidestar[]{%
      c_0\neq c_1\land(c_0\neq c_2\land c_1\neq c_2)}{%
      \rAI{12,\rAI{21,30}}}

\closeproofpartwideindR


  \proofpartwide{Reservenausschluss und Trägerinklusion}
    \proofstepwidestar[]{0\in\mathbb N}{%
      \FormulaRefAuto{PeanoZero}[ax]}
    \proofstepwidestar[]{1\in\mathbb N}{%
      \FormulaRefAuto{PeanoOneInN}[thm]}
    \proofstepwidestar[]{%
      c_0,c_1,c_2\notin U}{%
      \FormulaRefAuto{B21ReserveAvoidsFirstTwoRows}[thm]{1,2},\ \FormulaRefAuto{MogiljanskajaDPrimeDef}[def]}
    \proofstepwidestar[]{%
      U\subseteq D}{%
      \FormulaRefAuto{MogiljanskajaDPrimeDef}[def]}
    \proofstepwidestar[]{%
      P\subseteq D\land\{c_2\}\subseteq D}{%
      \FormulaRefAuto{B21ReserveMarkerDistinct}[thm],\ \FormulaRefAuto{MogiljanskajaDPrimeDef}[def],\ \FormulaRefAuto{PairSetSubsetFromMembership}[thm]}
    \proofstepwidestar[]{%
      V\subseteq D}{%
      \FormulaRefAuto{UnionSubsetCommonSuperset}[thm]{4,\rAEb{5}},\ \FormulaRefAuto{MogiljanskajaDPrimeDef}[def]}
    \proofstepwidestar[]{%
      D'\subseteq D}{%
      \FormulaRefAuto{UnionSubsetCommonSuperset}[thm]{\rAEa{5},6},\ \FormulaRefAuto{MogiljanskajaDPrimeDef}[def]}
  \closeproofpartwide

  \Needspace{6\baselineskip}
  \proofpartwide{Das ausgesparte Rechteckelement}
    \proofstepwidestar[]{0 \in \mathbb N}{%
      \FormulaRefAuto{PeanoZero}[ax]}
    \proofstepwidestar[]{1 \in \mathbb N}{%
      \FormulaRefAuto{PeanoOneInN}[thm]}
    \proofstepwidestar[]{%
      b_{01}\notin U}{%
      \FormulaRefAuto{B21ReserveAvoidsFirstTwoRows}[thm]{1,2}}
    \proofstepwidestar[4]{%
      b_{01}=c_0}{%
      \rA}
    \proofstepwidestar[]{0\in\mathbb N}{%
      \FormulaRefAuto{PeanoZero}[ax]}
    \proofstepwidestar[]{1\in\mathbb N}{%
      \FormulaRefAuto{PeanoOneInN}[thm]}
    \proofstepwidestar[4]{b_{01}=b_{00}}{%
      \FormulaRefAuto{MogiljanskajaDPrimeDef}[def]{4}}
    \proofstepwidestar[4]{0=0\land1=0}{%
      \FormulaRefAuto{MogiljanskajaBIndexUnique}[thm]{5,6,7}}
    \proofstepwidestar[4]{%
      1=0}{\rAEb{8}
}
    \proofstepwidestar[4]{1 \neq 0}{%
      \FormulaRefAuto{PeanoOneNeZero}[thm]}
    \proofstepwidestar[4]{%
      \bot}{%
      \rBI{10,9}}
    \proofstepwidestar[]{%
      b_{01}\neq c_0}{%
      \rCI{4,11}}
    \proofstepwidestar[13]{%
      b_{01}=c_1}{%
      \rA}
    \proofstepwidestar[]{0\in\mathbb N}{%
      \FormulaRefAuto{PeanoZero}[ax]}
    \proofstepwidestar[]{1\in\mathbb N}{%
      \FormulaRefAuto{PeanoOneInN}[thm]}
    \proofstepwidestar[13]{b_{01}=b_{11}}{%
      \FormulaRefAuto{MogiljanskajaDPrimeDef}[def]{13}}
    \proofstepwidestar[13]{0=1\land1=1}{%
      \FormulaRefAuto{MogiljanskajaBIndexUnique}[thm]{14,15,16}}
    \proofstepwidestar[13]{%
      0=1}{\rAEa{17}
}
    \proofstepwidestar[13]{0 \neq 1}{%
      \FormulaRefAuto{PeanoZeroNeOne}[thm]}
    \proofstepwidestar[13]{%
      \bot}{%
      \rBI{19,18}}
    \proofstepwidestar[]{%
      b_{01}\neq c_1}{%
      \rCI{13,20}}
    \proofstepwidestar[22]{%
      b_{01}=c_2}{%
      \rA}
    \proofstepwidestar[]{0\in\mathbb N}{%
      \FormulaRefAuto{PeanoZero}[ax]}
    \proofstepwidestar[]{1\in\mathbb N}{%
      \FormulaRefAuto{PeanoOneInN}[thm]}
    \proofstepwidestar[22]{b_{01}=b_{10}}{%
      \FormulaRefAuto{MogiljanskajaDPrimeDef}[def]{22}}
    \proofstepwidestar[22]{0=1\land1=0}{%
      \FormulaRefAuto{MogiljanskajaBIndexUnique}[thm]{23,24,25}}
    \proofstepwidestar[22]{%
      0=1}{\rAEa{26}
}
    \proofstepwidestar[22]{0 \neq 1}{%
      \FormulaRefAuto{PeanoZeroNeOne}[thm]}
    \proofstepwidestar[22]{%
      \bot}{%
      \rBI{28,27}}
    \proofstepwidestar[]{%
      b_{01}\neq c_2}{%
      \rCI{22,29}}
    \proofstepwidestar[]{%
      b_{01}\notin P\land b_{01}\notin V}{%
      \FormulaRefAuto{NotInPairSet}[thm]{12,21},\ \FormulaRefAuto{NotInUnion}[thm]{3,30},\ \FormulaRefAuto{MogiljanskajaDPrimeDef}[def]}
    \proofstepwidestar[]{%
      b_{01}\notin D'}{%
      \FormulaRefAuto{NotInUnion}[thm]{\rAEa{31},\rAEb{31}},\ \FormulaRefAuto{MogiljanskajaDPrimeDef}[def]}
  \closeproofpartwide
  \enlargethispage{2\baselineskip}\nopagebreak
\end{tabproofsplitwide}

\MogStatementTarget{MogiljanskajaReserveCoreFacts}
\FormulaThmDeltaKR[Der dritte Marker liegt außerhalb des festen Kerns]{%
  \begin{aligned}[t]
    &\text{(i)}\quad c_2\notin P,\\[-2pt]
    &\text{(ii)}\quad P\cap U=\varnothing.
  \end{aligned}%
}{%
  c_2\notin P\dsep P\cap U=\varnothing%
}{MogiljanskajaReserveCoreFacts}{%
  \DeltaRow{Mengen}{P\dsep U\dsep c_0\dsep c_1\dsep c_2}
}
\begin{tabproofsplitwide}
  \proofpartwide{Kern und Reserve sind disjunkt}
    \proofstepwidestar[]{%
      c_2\neq c_0\land c_2\neq c_1}{%
      \FormulaRefAuto{B21ReserveMarkerDistinct}[thm],\ \FormulaRefAuto{a=b\vdash b=a}[thm]}
    \proofstepwidestar[]{%
      c_2\notin P}{%
      \FormulaRefAuto{NotInPairSet}[thm]{1},\ \FormulaRefAuto{MogiljanskajaDPrimeDef}[def]}
    \proofstepwidestar[]{%
      c_0\notin U\land c_1\notin U}{%
      \FormulaRefAuto{MogiljanskajaReserveMarkerFacts}[thm]}
    \proofstepwidestar[]{%
      U\cap\{c_0,c_1\}=\varnothing}{%
      \FormulaRefAuto{PairSetDisjointFromTwoAvoided}[thm]{3}}
    \proofstepwidestar[]{%
      P\cap U=\varnothing}{%
      \FormulaRefAuto{MogiljanskajaDPrimeDef}[def]{\rIE{\FormulaRefAuto{A\cap B=B\cap A}[thm],4}}}
  \closeproofpartwide
\end{tabproofsplitwide}

\Needspace{20\baselineskip}
\section{Parametrisierungen und ihre Bildmengen}

\MogStatementTarget{MogiljanskajaRowParametrizationsDef}
\FormulaDefDeltaK[Die beiden Parametrisierungen der Produktschichten]{%
  \begin{aligned}[t]
    &\delta\colon\mathbb N\times\mathbb N\to D,\\[-2pt]
    &\forall n,j\in\mathbb N\,
      \bigl(\delta(n,j)\coloneqq b_{nj}\bigr);\\[2pt]
    &\upsilon\colon\mathbb N\times\mathbb N\to U,\\[-2pt]
    &\forall n,j\in\mathbb N\,
      \bigl(\upsilon(n,j)\coloneqq b_{(n+1)+1,j}\bigr)
  \end{aligned}%
}{MogiljanskajaRowParametrizationsDef}{%
  \DeltaRow{Mengen}{D\dsep U\dsep b_{ij}\dsep n\dsep j}
  \DeltaRow{Funktionen}{b\dsep\delta\dsep\upsilon}
}

\MogStatementTarget{MogiljanskajaRowParametrizationsBijective}
\hypertarget{mog.statement.B21DeltaBijective}{}
\hypertarget{mog.statement.B21UpsilonBijective}{}
\FormulaThmDeltaKR[Die Zeilenparametrisierungen sind bijektiv]{%
  \begin{aligned}[t]
    &\text{(i)}\quad \delta\colon\mathbb N\times\mathbb N\bij D,\\[-2pt]
    &\text{(ii)}\quad \upsilon\colon\mathbb N\times\mathbb N\bij U.
  \end{aligned}%
}{%
  \delta\colon\mathbb N\times\mathbb N\bij D\dsep
  \upsilon\colon\mathbb N\times\mathbb N\bij U%
}{MogiljanskajaRowParametrizationsBijective}{%
  \DeltaRow{Mengen}{D\dsep U\dsep n\dsep j\dsep m\dsep q\dsep u}
  \DeltaRow{Funktionen}{\delta\dsep\upsilon}
}
\begin{tabproofsplitwide}
  \proofpartwideindR[Unverschobene Zeilen]{%
    \delta\colon\mathbb N\times\mathbb N\bij D
  }{\delta\colon\mathbb N\times\mathbb N\bij D}[B21DeltaBijective]
    \proofstepwidestar[]{%
      \delta\colon\mathbb N\times\mathbb N\to D}{%
      \FormulaRefAuto{MogiljanskajaRowParametrizationsDef}[def]}
    \proofstepwidestar[2]{%
      n,j,m,q\in\mathbb N}{%
      \rA}
    \proofstepwidestar[3]{%
      \delta(n,j)=\delta(m,q)}{%
      \rA}
    \proofstepwidestar[2,3]{%
      b_{n,j}=b_{m,q}}{%
      \FormulaRefAuto{MogiljanskajaRowParametrizationsDef}[def]{2,3}}
    \proofstepwidestar[2,3]{%
      n=m\land j=q}{\FormulaRefAuto{MogiljanskajaBIndexUnique}[thm]{2,4}
}
    \proofstepwidestar[2,3]{%
      n=m\land j=q}{%
      \rRE{\FormulaRefAuto{P\rightarrow P}[thm],5}}
    \proofstepwidestar[2,3]{%
      (n,j)=(m,q)}{%
      \FormulaRefAuto{(a,b)=(c,d)\eqvdash a=c\land b=d}[thm]{6}}
    \proofstepwidestar[]{%
      \delta\colon\mathbb N\times\mathbb N\inj D}{%
      \FormulaRefAuto{Injektive Funktion}[def]{1,\rUI{\rUI{\rUI{\rUI{\rRI{2,\rRI{3,7}}}}}}}}
    \proofstepwidestar[9]{%
      u\in D}{%
      \rA}
    \proofstepwidestar[9]{%
      \exists n,j\in\mathbb N\,(u=b_{n,j})}{%
      \FormulaRefAuto{MogiljanskajaLayerDMembership}[thm]{9}}
    \proofstepwidestar[11]{%
      n\in\mathbb N\land(j\in\mathbb N\land u=b_{n,j})}{%
      \rA}
    \proofstepwidestar[11]{%
      (n,j)\in\mathbb N\times\mathbb N\land \delta(n,j)=u}{%
      \FormulaRefAuto{MogiljanskajaRowParametrizationsDef}[def]{11},\ \FormulaRefAuto{x\in A\times B\eqvdash\exists a\in A\,\exists b\in B\,x=(a,b)}[thm],\ \FormulaRefAuto{a=b\vdash b=a}[thm]}
    \proofstepwidestar[11]{%
      \exists p\in\mathbb N\times\mathbb N\,(\delta(p)=u)}{%
      \rEI{12}}
    \proofstepwidestar[9]{%
      \exists p\in\mathbb N\times\mathbb N\,(\delta(p)=u)}{%
      \rEE{10,11,13}}
    \proofstepwidestar[]{%
      \delta\colon\mathbb N\times\mathbb N\sur D}{%
      \FormulaRefAuto{Surjektive Funktion}[def]{1,\rUI{\rRI{9,14}}}}
    \proofstepwidestar[]{%
      \delta\colon\mathbb N\times\mathbb N\bij D}{%
      \FormulaRefAuto{Bijektive Funktion}[def]{8,15}}
  \closeproofpartwideindR

  \proofpartwideindR[Verschobene Zeilen]{%
    \upsilon\colon\mathbb N\times\mathbb N\bij U
  }{\upsilon\colon\mathbb N\times\mathbb N\bij U}[B21UpsilonBijective]
    \proofstepwidestar[]{%
      \upsilon\colon\mathbb N\times\mathbb N\to U}{%
      \FormulaRefAuto{MogiljanskajaRowParametrizationsDef}[def]}
    \proofstepwidestar[2]{%
      n,j,m,q\in\mathbb N}{%
      \rA}
    \proofstepwidestar[3]{%
      \upsilon(n,j)=\upsilon(m,q)}{%
      \rA}
    \proofstepwidestar[2,3]{%
      b_{(n+1)+1,j}=b_{(m+1)+1,q}}{%
      \FormulaRefAuto{MogiljanskajaRowParametrizationsDef}[def]{2,3}}
    \proofstepwidestar[2]{n+1\in\mathbb N}{%
      \FormulaRefAuto{PeanoAddOneClosure}[thm]{\rAEa{2}}}
    \proofstepwidestar[2]{m+1\in\mathbb N}{%
      \FormulaRefAuto{PeanoAddOneClosure}[thm]{\rAEa{\rAEb{\rAEb{2}}}}}
    \proofstepwidestar[2]{(n+1)+1\in\mathbb N}{%
      \FormulaRefAuto{PeanoAddOneClosure}[thm]{5}}
    \proofstepwidestar[2]{(m+1)+1\in\mathbb N}{%
      \FormulaRefAuto{PeanoAddOneClosure}[thm]{6}}
    \proofstepwidestar[2,3]{%
      (n+1)+1=(m+1)+1\land j=q}{\FormulaRefAuto{MogiljanskajaBIndexUnique}[thm]{2,4,\rAI{7,8}}
}
    \proofstepwidestar[2,3]{n+1=m+1}{%
      \FormulaRefAuto{PeanoAddOneInjective}[thm]{\rAEa{9}}}
    \proofstepwidestar[2,3]{%
      n=m\land j=q}{%
      \FormulaRefAuto{PeanoAddOneInjective}[thm]{10},\ \rAEb{9}}
    \proofstepwidestar[2,3]{%
      (n,j)=(m,q)}{%
      \FormulaRefAuto{(a,b)=(c,d)\eqvdash a=c\land b=d}[thm]{11}}
    \proofstepwidestar[]{%
      \upsilon\colon\mathbb N\times\mathbb N\inj U}{%
      \FormulaRefAuto{Injektive Funktion}[def]{1,\rUI{\rUI{\rUI{\rUI{\rRI{2,\rRI{3,12}}}}}}}}
    \proofstepwidestar[14]{%
      u\in U}{%
      \rA}
    \proofstepwidestar[14]{%
      \exists n,j\in\mathbb N\,(u=b_{(n+1)+1,j})}{%
      \FormulaRefAuto{MogiljanskajaReserveMembership}[thm]{14}}
    \proofstepwidestar[16]{%
      n\in\mathbb N\land(j\in\mathbb N\land u=b_{(n+1)+1,j})}{%
      \rA}
    \proofstepwidestar[16]{%
      (n,j)\in\mathbb N\times\mathbb N\land \upsilon(n,j)=u}{%
      \FormulaRefAuto{MogiljanskajaRowParametrizationsDef}[def]{16},\ \FormulaRefAuto{x\in A\times B\eqvdash\exists a\in A\,\exists b\in B\,x=(a,b)}[thm],\ \FormulaRefAuto{a=b\vdash b=a}[thm]}
    \proofstepwidestar[16]{%
      \exists p\in\mathbb N\times\mathbb N\,(\upsilon(p)=u)}{%
      \rEI{17}}
    \proofstepwidestar[14]{%
      \exists p\in\mathbb N\times\mathbb N\,(\upsilon(p)=u)}{%
      \rEE{15,16,18}}
    \proofstepwidestar[]{%
      \upsilon\colon\mathbb N\times\mathbb N\sur U}{%
      \FormulaRefAuto{Surjektive Funktion}[def]{1,\rUI{\rRI{14,19}}}}
    \proofstepwidestar[]{%
      \upsilon\colon\mathbb N\times\mathbb N\bij U}{%
      \FormulaRefAuto{Bijektive Funktion}[def]{13,20}}
  \closeproofpartwideindR
\end{tabproofsplitwide}

Insbesondere gelten die präzisen Bildmengengleichheiten
\[
  \delta[\mathbb N\times\mathbb N]=D,
  \qquad
  \upsilon[\mathbb N\times\mathbb N]=U.
\]

\MogStatementTarget{MogiljanskajaThetaDef}
\FormulaDefDeltaK[Die Zeilenrückverschiebung]{%
  \theta\coloneqq\delta\circ\upsilon^{-1}%
}{MogiljanskajaThetaDef}{%
  \DeltaRow{Mengen}{D\dsep U}
  \DeltaRow{Funktionen}{\delta\dsep\upsilon\dsep\theta}
}

\MogStatementTarget{MogiljanskajaThetaBijective}
\FormulaThmDeltaK[Die Zeilenrückverschiebung ist bijektiv]{%
  \theta\colon U\bij D%
}{MogiljanskajaThetaBijective}{%
  \DeltaRow{Mengen}{D\dsep U}
  \DeltaRow{Funktionen}{\delta\dsep\upsilon\dsep\theta}
}
\begin{tabproofsplitwide}
  \proofpartwide{Inverse und Komposition}
    \proofstepwidestar[]{%
      \upsilon\colon\mathbb N\times\mathbb N\bij U}{%
      \FormulaRefAuto{B21UpsilonBijective}[thm]}
    \proofstepwidestar[]{%
      \upsilon^{-1}\colon U\bij\mathbb N\times\mathbb N}{%
      \FormulaRefAuto{InverseFunctionBijection}[thm]{1}}
    \proofstepwidestar[]{%
      \delta\colon\mathbb N\times\mathbb N\bij D}{%
      \FormulaRefAuto{B21DeltaBijective}[thm]}
    \proofstepwidestar[]{%
      \delta\circ\upsilon^{-1}\colon U\bij D}{%
      \FormulaRefAuto{BijectiveFunctionComposition}[thm]{2,3}}
    \proofstepwidestar[]{%
      \theta\colon U\bij D}{%
      \FormulaRefAuto{MogiljanskajaThetaDef}[def],\ \rIE{4}}
  \closeproofpartwide
\end{tabproofsplitwide}

\Needspace{20\baselineskip}
\section{Die punktierte Verschiebung}

\MogStatementTarget{MogiljanskajaShiftValueInV}
\FormulaThmDeltaK[Die Werte der Verschiebungsfolge liegen in der Reserve]{%
  n\in\mathbb N\vdash b_{n+1,0}\in V%
}{MogiljanskajaShiftValueInV}{%
  \DeltaRow{Mengen}{U\dsep V\dsep c_2\dsep b_{ij}\dsep n}
  \DeltaRow{Funktionen}{b}
}
\begin{tabproofsplitwide}
  \proofpartwide{Nullindex und Nachfolgerindex}
    \proofstepwidestar[1]{%
      n\in\mathbb N}{%
      \rA}
    \proofstepwidestar[1]{%
      n=0\lor\exists m\in\mathbb N\,(n=m+1)}{%
      \FormulaRefAuto{PeanoIndexSplitAtOne}[thm]{1}}
    \proofstepwidestar[3]{%
      n=0}{%
      \rA}
    \proofstepwidestar[]{1\in\mathbb N}{%
      \FormulaRefAuto{PeanoOneInN}[thm]}
    \proofstepwidestar[3]{%
      b_{n+1,0}=c_2}{%
      \rIE{3},\ \FormulaRefAuto{PeanoAddLeftZero}[thm]{4},\ \FormulaRefAuto{MogiljanskajaDPrimeDef}[def]}
    \proofstepwidestar[3]{%
      b_{n+1,0}\in V}{%
      \FormulaRefAuto{MogiljanskajaDPrimeDef}[def]{5},\ \FormulaRefAuto{x\in A\cup B\eqvdash x\in A\lor x\in B}[thm]}
    \proofstepwidestar[7]{%
      \exists m\in\mathbb N\,(n=m+1)}{%
      \rA}
    \proofstepwidestar[8]{%
      m\in\mathbb N\land n=m+1}{%
      \rA}
    \proofstepwidestar[8]{%
      n=m+1}{%
      \rAEb{8}}
    \proofstepwidestar[8]{0 \in \mathbb N}{%
      \FormulaRefAuto{PeanoZero}[ax]}
    \proofstepwidestar[8]{%
      b_{n+1,0}\in U}{%
      \FormulaRefAuto{MogiljanskajaReserveMembership}[thm]{\rEI{\rAI{\rAEa{8},\rEI{\rAI{10,\rIE{9,\rII}}}}}}}
    \proofstepwidestar[8]{%
      b_{n+1,0}\in V}{%
      \FormulaRefAuto{MogiljanskajaDPrimeDef}[def]{11},\ \FormulaRefAuto{x\in A\cup B\eqvdash x\in A\lor x\in B}[thm]}
    \proofstepwidestar[7]{%
      b_{n+1,0}\in V}{%
      \rEE{7,8,12}}
    \proofstepwidestar[1]{%
      b_{n+1,0}\in V}{%
      \rOE{2,3,6,7,13}}
  \closeproofpartwide
\end{tabproofsplitwide}

\MogStatementTarget{MogiljanskajaShiftSequenceDef}
\FormulaDefDeltaK[Die Folge für die punktierte Verschiebung]{%
  \begin{aligned}[t]
    &H\colon\mathbb N\to V,\\[-2pt]
    &\forall n\in\mathbb N\,
      \bigl(H_n\coloneqq b_{n+1,0}\bigr)
  \end{aligned}%
}{MogiljanskajaShiftSequenceDef}{%
  \DeltaRow{Mengen}{V\dsep b_{ij}\dsep n}
  \DeltaRow{Funktionen}{b\dsep H}
}

\MogStatementTarget{MogiljanskajaShiftSequenceFacts}
\hypertarget{mog.statement.B21ShiftSequenceInjective}{}
\FormulaThmDeltaKR[Die punktierte Verschiebung entfernt genau den dritten
Marker]{%
  \begin{aligned}[t]
    &\text{(i)}\quad H\colon\mathbb N\inj V,\\[-2pt]
    &\text{(ii)}\quad H_0=c_2,\\[-2pt]
    &\text{(iii)}\quad V\setminus\{c_2\}=U.
  \end{aligned}%
}{%
  H\colon\mathbb N\inj V\dsep H_0=c_2\dsep
  V\setminus\{c_2\}=U%
}{MogiljanskajaShiftSequenceFacts}{%
  \DeltaRow{Mengen}{U\dsep V\dsep c_2\dsep n\dsep m}
  \DeltaRow{Funktionen}{H}
}
\begin{tabproofsplitwide}
  \proofpartwideindR[Injektivität der Verschiebungsfolge]{%
    H\colon\mathbb N\inj V
  }{H\colon\mathbb N\inj V}[B21ShiftSequenceInjective]
    \proofstepwidestar[]{%
      H\colon\mathbb N\to V}{%
      \FormulaRefAuto{MogiljanskajaShiftSequenceDef}[def]}
    \proofstepwidestar[2]{%
      n,m\in\mathbb N}{%
      \rA}
    \proofstepwidestar[3]{%
      H_n=H_m}{%
      \rA}
    \proofstepwidestar[2,3]{%
      b_{n+1,0}=b_{m+1,0}}{%
      \FormulaRefAuto{MogiljanskajaShiftSequenceDef}[def]{2,3}}
    \proofstepwidestar[2,3]{0 \in \mathbb N}{%
      \FormulaRefAuto{PeanoZero}[ax]}
    \proofstepwidestar[2]{n+1\in\mathbb N}{%
      \FormulaRefAuto{PeanoAddOneClosure}[thm]{\rAEa{2}}}
    \proofstepwidestar[2]{m+1\in\mathbb N}{%
      \FormulaRefAuto{PeanoAddOneClosure}[thm]{\rAEb{2}}}
    \proofstepwidestar[2,3]{n+1=m+1\land0=0}{%
      \FormulaRefAuto{MogiljanskajaBIndexUnique}[thm]{2,\rAI{6,7},5,4}}
    \proofstepwidestar[2,3]{%
      n+1=m+1}{\rAEa{8}
}
    \proofstepwidestar[2,3]{%
      n=m}{%
      \FormulaRefAuto{PeanoAddOneInjective}[thm]{2,9}}
    \proofstepwidestar[]{%
      H\colon\mathbb N\inj V}{%
      \FormulaRefAuto{InjectiveIndexedFamilyCriterion}[thm]{1,\rUI{\rUI{\rRI{2,\rRI{3,10}}}}}}
  \closeproofpartwideindR

  \proofpartwide{Anfangswert und punktierte Zielmenge}
    \proofstepwidestar[]{0\in\mathbb N}{%
      \FormulaRefAuto{PeanoZero}[ax]}
    \proofstepwidestar[]{1\in\mathbb N}{%
      \FormulaRefAuto{PeanoOneInN}[thm]}
    \proofstepwidestar[]{%
      H_0=b_{0+1,0}=c_2}{%
      \FormulaRefAuto{MogiljanskajaShiftSequenceDef}[def]{1},\ \FormulaRefAuto{PeanoAddLeftZero}[thm]{2},\ \FormulaRefAuto{MogiljanskajaDPrimeDef}[def]}
    \proofstepwidestar[]{%
      c_2\notin U}{%
      \FormulaRefAuto{MogiljanskajaReserveMarkerFacts}[thm]}
    \proofstepwidestar[]{%
      \{c_2\}\cap U=\varnothing}{%
      \FormulaRefAuto{SingletonDisjointFromAvoidedSet}[thm]{4}}
    \proofstepwidestar[]{%
      (\{c_2\}\cup U)\setminus\{c_2\}=U}{%
      \FormulaRefAuto{RemoveDisjointUnionSummand}[thm]{5}}
    \proofstepwidestar[]{%
      V\setminus\{c_2\}=U}{%
      \FormulaRefAuto{MogiljanskajaDPrimeDef}[def]{\rIE{\FormulaRefAuto{A\cup B=B\cup A}[thm],6}}}
  \closeproofpartwide
\end{tabproofsplitwide}

\MogStatementTarget{B21ShiftImageTermSetDef}
\FormulaDefDeltaKR[Die durch die Verschiebungswerte beschriebene Menge]{%
  \begin{aligned}[t]
    &\{b_{n+1,0}\mid n\in\mathbb N\}\coloneqq\\[-2pt]
    &\quad\iota S\,\forall x\,
      \bigl(x\in S\leftrightarrow\exists n\in\mathbb N\,(x=b_{n+1,0})\bigr)
  \end{aligned}%
}{%
  \{b_{n+1,0}\mid n\in\mathbb N\}\coloneqq
  \iota S\,\forall x\,
    \bigl(x\in S\leftrightarrow\exists n\in\mathbb N\,(x=b_{n+1,0})\bigr)%
}{B21ShiftImageTermSetDef}{%
  \DeltaRow{Mengen}{S\dsep x\dsep n\dsep b_{ij}}
  \DeltaRow{Eindeutige Existenz}{%
    \begin{aligned}[t]
      &\exists!S\,\forall x\\[-2pt]
      &\quad\bigl(x\in S\leftrightarrow\exists n\in\mathbb N\,(x=b_{n+1,0})\bigr)
    \end{aligned}}
    [\FormulaRefAuto{TermImageSetUniqueExists}[thm]]
}

\MogStatementTarget{MogiljanskajaShiftSequenceImage}
\FormulaThmDeltaKR[Tatsächliche Bildmenge der Verschiebungsfolge]{%
  \begin{aligned}[t]
    &\text{(i)}\quad H[\mathbb N]=\{b_{n+1,0}\mid n\in\mathbb N\},\\[-2pt]
    &\text{(ii)}\quad H[\mathbb N]\subsetneq V.
  \end{aligned}%
}{%
  H[\mathbb N]
  =\{b_{n+1,0}\mid n\in\mathbb N\}
  \dsep H[\mathbb N]\subsetneq V%
}{MogiljanskajaShiftSequenceImage}{%
  \DeltaRow{Mengen}{V\dsep b_{ij}\dsep n}
  \DeltaRow{Funktionen}{H}
}
\begin{tabproofsplitwide}
  \proofpartwide{Bildgleichheit und Inklusion}
    \proofstepwidestar[]{%
      H\colon\mathbb N\to V}{%
      \FormulaRefAuto{MogiljanskajaShiftSequenceDef}[def]}
    \proofstepwidestar[]{%
      x\in H[\mathbb N]\leftrightarrow\exists n\in\mathbb N\,(x=b_{n+1,0})}{%
      \FormulaRefAuto{IndexedFamilyImageMembership}[thm]{1},\ \FormulaRefAuto{IndexedFamilyImageDef}[def],\ \FormulaRefAuto{MogiljanskajaShiftSequenceDef}[def]}
    \proofstepwidestar[]{%
      H[\mathbb N]=\{b_{n+1,0}\mid n\in\mathbb N\}}{%
      \FormulaRefAuto{\forall x\,(x\in A\leftrightarrow x\in B)\vdash A=B}[ax]{\rUI{\rLRS{\FormulaRefAuto{B21ShiftImageTermSetDef}[def],2}}}}
    \proofstepwidestar[]{%
      H[\mathbb N]\subseteq V}{%
      \FormulaRefAuto{F\colon A\to B\vdash F[A]\subseteq B}[thm]{1}}
  \closeproofpartwide

  \proofpartwide{Ein explizit ausgelassener Wert}
    \proofstepwidestar[]{0 \in \mathbb N}{%
      \FormulaRefAuto{PeanoZero}[ax]}
    \proofstepwidestar[]{1 \in \mathbb N}{%
      \FormulaRefAuto{PeanoOneInN}[thm]}
    \proofstepwidestar[]{%
      b_{(0+1)+1,1}\in U}{%
      \FormulaRefAuto{MogiljanskajaReserveMembership}[thm]{\rEI{\rAI{1,\rEI{\rAI{2,\rII}}}}}}
    \proofstepwidestar[]{%
      b_{(0+1)+1,1}\in V}{%
      \FormulaRefAuto{MogiljanskajaDPrimeDef}[def]{3},\ \FormulaRefAuto{x\in A\cup B\eqvdash x\in A\lor x\in B}[thm]}
    \proofstepwidestar[5]{%
      b_{(0+1)+1,1}\in H[\mathbb N]}{%
      \rA}
    \proofstepwidestar[]{H\colon\mathbb N\to V}{%
      \FormulaRefAuto{MogiljanskajaShiftSequenceDef}[def]}
    \proofstepwidestar[5]{%
      \exists n\in\mathbb N\,(b_{(0+1)+1,1}=H_n)}{%
      \FormulaRefAuto{IndexedFamilyImageMembership}[thm]{6,5},\ \FormulaRefAuto{IndexedFamilyImageDef}[def]}
    \proofstepwidestar[8]{%
      n\in\mathbb N\land b_{(0+1)+1,1}=H_n}{%
      \rA}
    \proofstepwidestar[8]{%
      b_{(0+1)+1,1}=b_{n+1,0}}{%
      \FormulaRefAuto{MogiljanskajaShiftSequenceDef}[def]{\rAEa{8}},\ \rIE{\rAEb{8}}}
    \proofstepwidestar[8]{0 \in \mathbb N}{%
      \FormulaRefAuto{PeanoZero}[ax]}
    \proofstepwidestar[8]{1 \in \mathbb N}{%
      \FormulaRefAuto{PeanoOneInN}[thm]}
    \proofstepwidestar[8]{0+1\in\mathbb N}{%
      \FormulaRefAuto{PeanoAddOneClosure}[thm]{10}}
    \proofstepwidestar[8]{n+1\in\mathbb N}{%
      \FormulaRefAuto{PeanoAddOneClosure}[thm]{\rAEa{8}}}
    \proofstepwidestar[8]{(0+1)+1\in\mathbb N}{%
      \FormulaRefAuto{PeanoAddOneClosure}[thm]{12}}
    \proofstepwidestar[8]{(0+1)+1=n+1\land1=0}{%
      \FormulaRefAuto{MogiljanskajaBIndexUnique}[thm]{10,11,\rAI{14,13},8,9}}
    \proofstepwidestar[8]{%
      1=0}{\rAEb{15}
}
    \proofstepwidestar[8]{1 \neq 0}{%
      \FormulaRefAuto{PeanoOneNeZero}[thm]}
    \proofstepwidestar[8]{%
      \bot}{%
      \rBI{17,16}}
    \proofstepwidestar[5]{%
      \bot}{%
      \rEE{7,8,18}}
    \proofstepwidestar[]{%
      b_{(0+1)+1,1}\notin H[\mathbb N]}{%
      \rCI{5,19}}
    \proofstepwidestar[]{H\colon\mathbb N\to V}{%
      \FormulaRefAuto{MogiljanskajaShiftSequenceDef}[def]}
    \proofstepwidestar[]{%
      H[\mathbb N]\subseteq V}{%
      \FormulaRefAuto{IndexedFamilyImageSubsetCodomain}[thm]{21},\ \FormulaRefAuto{IndexedFamilyImageDef}[def]}
    \proofstepwidestar[]{%
      H[\mathbb N]\subsetneq V}{%
      \FormulaRefAuto{A\subset B\coloneq A\subseteq B\land A\neq B}[def]{\rAI{22,\FormulaRefAuto{x\in B, x\not\in A\vdash A\neq B}[thm]{4,20}}}}
  \closeproofpartwide
\end{tabproofsplitwide}

Damit ist ausdrücklich zwischen Ziel- und Bildmenge unterschieden:
\(H\colon\mathbb N\to V\), aber \(H[\mathbb N]\neq V\). Die Folge wird
nicht zur Aufzählung von \(V\), sondern zur Konstruktion einer Verschiebung
auf \(V\) verwendet.

\MogStatementTarget{MogiljanskajaSigmaDef}
\FormulaDefDeltaK[Die Bijektion in die um einen Marker erweiterte Reserve]{%
  \sigma\coloneqq(\SuccShift{H})^{-1}%
}{MogiljanskajaSigmaDef}{%
  \DeltaRow{Mengen}{U\dsep V}
  \DeltaRow{Funktionen}{H\dsep\SuccShift{H}\dsep\sigma}
}

\MogStatementTarget{MogiljanskajaSigmaBijective}
\FormulaThmDeltaK[Die Markeradjunktion ist bijektiv]{%
  \sigma\colon U\bij V%
}{MogiljanskajaSigmaBijective}{%
  \DeltaRow{Mengen}{U\dsep V\dsep c_2}
  \DeltaRow{Funktionen}{H\dsep\SuccShift{H}\dsep\sigma}
}
\begin{tabproofsplitwide}
  \proofpartwide{Inverse der punktierten Verschiebung}
    \proofstepwidestar[]{%
      H\colon\mathbb N\inj V}{%
      \FormulaRefAuto{B21ShiftSequenceInjective}[thm]}
    \proofstepwidestar[]{%
      H_0=c_2\land V\setminus\{c_2\}=U}{%
      \FormulaRefAuto{MogiljanskajaShiftSequenceFacts}[thm]}
    \proofstepwidestar[]{%
      \SuccShift{H}\colon V\bij V\setminus\{H(0)\}}{%
      \FormulaRefAuto{SuccShiftPuncturedBijection}[thm]{1}}
    \proofstepwidestar[]{0\in\mathbb N}{%
      \FormulaRefAuto{PeanoZero}[ax]}
    \proofstepwidestar[]{H\colon\mathbb N\to V}{%
      \FormulaRefAuto{MogiljanskajaShiftSequenceDef}[def]}
    \proofstepwidestar[]{%
      \SuccShift{H}\colon V\bij U}{%
      \FormulaRefAuto{IndexedFamilyDef}[def]{5,4},\ \rIE{2,3}}
    \proofstepwidestar[]{%
      (\SuccShift{H})^{-1}\colon U\bij V}{%
      \FormulaRefAuto{InverseFunctionBijection}[thm]{6}}
    \proofstepwidestar[]{%
      \sigma\colon U\bij V}{%
      \FormulaRefAuto{MogiljanskajaSigmaDef}[def],\ \rIE{7}}
  \closeproofpartwide
\end{tabproofsplitwide}

\Needspace{20\baselineskip}
\section{Familien mit den vorgegebenen Bildmengen}

Die Familien \(a\), \(b\), \(\delta\) und \(\upsilon\) parametrisieren
bereits \(L,D,D\) beziehungsweise \(U\). Aus \(\sigma\) entsteht ohne eine
weitere Fallunterscheidung eine bijektive Familie mit Bild \(V\).

\MogStatementTarget{MogiljanskajaVParametrizationDef}
\FormulaDefDeltaK[Eine bijektive Familie mit Bildmenge (V)]{%
  \nu\coloneqq\sigma\circ\upsilon%
}{MogiljanskajaVParametrizationDef}{%
  \DeltaRow{Mengen}{U\dsep V}
  \DeltaRow{Funktionen}{\upsilon\dsep\sigma\dsep\nu}
}

\MogStatementTarget{MogiljanskajaVParametrizationBijective}
\FormulaThmDeltaKR[Die Familie (nu) parametrisiert (V)]{%
  \begin{aligned}[t]
    &\text{(i)}\quad \nu\colon\mathbb N\times\mathbb N\bij V,\\[-2pt]
    &\text{(ii)}\quad \nu[\mathbb N\times\mathbb N]=V.
  \end{aligned}%
}{%
  \nu\colon\mathbb N\times\mathbb N\bij V\dsep
  \nu[\mathbb N\times\mathbb N]=V%
}{MogiljanskajaVParametrizationBijective}{%
  \DeltaRow{Mengen}{V}
  \DeltaRow{Funktionen}{\nu}
}
\begin{tabproofsplitwide}
  \proofpartwide{Komposition und Bildmenge}
    \proofstepwidestar[]{%
      \upsilon\colon\mathbb N\times\mathbb N\bij U}{%
      \FormulaRefAuto{B21UpsilonBijective}[thm]}
    \proofstepwidestar[]{%
      \sigma\colon U\bij V}{%
      \FormulaRefAuto{MogiljanskajaSigmaBijective}[thm]}
    \proofstepwidestar[]{%
      \sigma\circ\upsilon\colon\mathbb N\times\mathbb N\bij V}{%
      \FormulaRefAuto{BijectiveFunctionComposition}[thm]{1,2}}
    \proofstepwidestar[]{%
      \nu\colon\mathbb N\times\mathbb N\bij V}{%
      \FormulaRefAuto{MogiljanskajaVParametrizationDef}[def],\ \rIE{3}}
    \proofstepwidestar[]{%
      \nu\colon\mathbb N\times\mathbb N\sur V}{%
      \FormulaRefAuto{Bijektive Funktion}[def]{4}}
    \proofstepwidestar[]{%
      \nu[\mathbb N\times\mathbb N]=V}{%
      \FormulaRefAuto{SurjectiveImageEqualsCodomain}[thm]{5}}
  \closeproofpartwide
\end{tabproofsplitwide}

Für \(D'=P\cup V\) müssen die beiden Marker aus \(P\) und die Familie
\(\nu\) disjunkt zusammengefügt werden. Die Tags \(0\) und \(1\) sorgen
dafür, dass sich die beiden Indexteile nicht überschneiden.

\MogStatementTarget{MogiljanskajaDPrimeIndexSetDef}
\FormulaDefDeltaK[Indexmenge für die Familie mit Bild (D')]{%
  I_{D'}\coloneqq
  (\{0\}\times\NatSeg{2})
  \cup
  (\{1\}\times(\mathbb N\times\mathbb N))%
}{MogiljanskajaDPrimeIndexSetDef}{%
  \DeltaRow{Mengen}{I_{D'}}
}

\MogStatementTarget{MogiljanskajaDPrimeParametrizationDef}
\FormulaDefDeltaK[Eine Familie mit Bildmenge (D')]{%
  \begin{aligned}[t]
    &\omega\colon I_{D'}\to D',\\[-2pt]
    &\forall i\in\NatSeg{2}\,
      \left(
        \omega((0,i))\coloneqq
        \begin{cases}
          c_0,& i=0,\\
          c_1,& i=1,
        \end{cases}
      \right),\\[-2pt]
    &\forall n,j\in\mathbb N\,
      \bigl(\omega((1,(n,j)))\coloneqq\nu(n,j)\bigr)
  \end{aligned}%
}{MogiljanskajaDPrimeParametrizationDef}{%
  \DeltaRow{Mengen}{I_{D'}\dsep D'\dsep c_0\dsep c_1\dsep n\dsep j}
  \DeltaRow{Funktionen}{\nu\dsep\omega}
}

\MogStatementTarget{B21TaggedTwoLayerEquality}
\FormulaThmDeltaK[Die mit null markierte Zweierschicht]{%
  \{0\}\times\NatSeg2=\{(0,0),(0,1)\}%
}{B21TaggedTwoLayerEquality}{%
  \DeltaRow{Mengen}{z\dsep i}
}
\begin{tabproofwide}
  \proofstepwidestar[]{z\in\{0\}\times\NatSeg2\leftrightarrow
    \exists i\in\NatSeg2\,(z=(0,i))}{%
    \FormulaRefAuto{TaggedProductLayerRepresentation}[thm]}
  \proofstepwidestar[2]{z\in\{0\}\times\NatSeg2}{\rA}
  \proofstepwidestar[2]{\exists i\in\NatSeg2\,(z=(0,i))}{\rRE{\rLREa{1},2}}
  \proofstepwidestar[4]{i\in\NatSeg2\land z=(0,i)}{\rA}
  \proofstepwidestar[4]{i=0\lor i=1}{%
    \FormulaRefAuto{PeanoNatSegTwoMembership}[thm]{\rAEa{4}}}
  \proofstepwidestar[6]{i=0}{\rA}
  \proofstepwidestar[4,6]{z=(0,0)}{\rIE{6,\rAEb{4}}}
  \proofstepwidestar[4,6]{z\in\{(0,0),(0,1)\}}{%
    \FormulaRefAuto{x\in\{a,b\}\eqvdash(x=a\lor x=b)}[thm]{\rOIa{7}}}
  \proofstepwidestar[9]{i=1}{\rA}
  \proofstepwidestar[4,9]{z=(0,1)}{\rIE{9,\rAEb{4}}}
  \proofstepwidestar[4,9]{z\in\{(0,0),(0,1)\}}{%
    \FormulaRefAuto{x\in\{a,b\}\eqvdash(x=a\lor x=b)}[thm]{\rOIb{10}}}
  \proofstepwidestar[4]{z\in\{(0,0),(0,1)\}}{\rOE{5,6,8,9,11}}
  \proofstepwidestar[2]{z\in\{(0,0),(0,1)\}}{\rEE{3,4,12}}
  \proofstepwidestar[14]{z\in\{(0,0),(0,1)\}}{\rA}
  \proofstepwidestar[14]{z=(0,0)\lor z=(0,1)}{%
    \FormulaRefAuto{x\in\{a,b\}\eqvdash(x=a\lor x=b)}[thm]{14}}
  \proofstepwidestar[16]{z=(0,0)}{\rA}
  \proofstepwidestar[16]{\exists i\in\NatSeg2\,(z=(0,i))}{%
    \rEI{\rAI{\FormulaRefAuto{PeanoNatSegTwoMembership}[thm]{\rOIa{\rII}},16}}}
  \proofstepwidestar[18]{z=(0,1)}{\rA}
  \proofstepwidestar[18]{\exists i\in\NatSeg2\,(z=(0,i))}{%
    \rEI{\rAI{\FormulaRefAuto{PeanoNatSegTwoMembership}[thm]{\rOIb{\rII}},18}}}
  \proofstepwidestar[14]{\exists i\in\NatSeg2\,(z=(0,i))}{\rOE{15,16,17,18,19}}
  \proofstepwidestar[14]{z\in\{0\}\times\NatSeg2}{\rRE{\rLREb{1},20}}
  \proofstepwidestar[]{z\in\{0\}\times\NatSeg2\leftrightarrow z\in\{(0,0),(0,1)\}}{%
    \rLRI{\rRI{2,13},\rRI{14,21}}}
  \proofstepwidestar[]{\{0\}\times\NatSeg2=\{(0,0),(0,1)\}}{%
    \FormulaRefAuto{\forall x\,(x\in A\leftrightarrow x\in B)\vdash A=B}[ax]{\rUI{22}}}
\end{tabproofwide}

\MogStatementTarget{MogiljanskajaDPrimeParametrizationBijective}
\hypertarget{mog.statement.B21DPrimeFamilyBijective}{}
\FormulaThmDeltaKR[Die Familie (omega) parametrisiert (D')]{%
  \begin{aligned}[t]
    &\text{(i)}\quad \omega\colon I_{D'}\bij D',\\[-2pt]
    &\text{(ii)}\quad \omega[I_{D'}]=D'.
  \end{aligned}%
}{%
  \omega\colon I_{D'}\bij D'\dsep \omega[I_{D'}]=D'%
}{MogiljanskajaDPrimeParametrizationBijective}{%
  \DeltaRow{Mengen}{I_{D'}\dsep D'\dsep P\dsep V}
  \DeltaRow{Funktionen}{\omega}
}
\begin{tabproofsplitwide}
  \proofpartwideindR[Die beiden Indexteile]{%
    \omega\colon I_{D'}\bij D'
  }{\omega\colon I_{D'}\bij D'}[B21DPrimeFamilyBijective]
    \proofstepwidestar[]{%
      I_0\coloneqq\{0\}\times\NatSeg 2}{%
      \text{Abkürzung}}
    \proofstepwidestar[]{%
      I_1\coloneqq\{1\}\times(\mathbb N\times\mathbb N)}{%
      \text{Abkürzung}}
    \proofstepwidestar[]{%
      I_{D'}=I_0\cup I_1}{%
      \FormulaRefAuto{MogiljanskajaDPrimeIndexSetDef}[def],\ 1,\ 2}
    \proofstepwidestar[]{0\neq1}{%
      \FormulaRefAuto{PeanoZeroNeOne}[thm]}
    \proofstepwidestar[]{%
      I_0\cap I_1=\varnothing}{%
      \FormulaRefAuto{TaggedProductLayersDisjoint}[thm]{4},\ 1,\ 2}
    \proofstepwidestar[]{%
      \omega\colon I_{D'}\to D'}{%
      \FormulaRefAuto{MogiljanskajaDPrimeParametrizationDef}[def]}
    \proofstepwidestar[]{2\in\mathbb N}{%
      \FormulaRefAuto{PeanoTwoInN}[thm]}
    \proofstepwidestar[]{%
      I_0=\{(0,0),(0,1)\}}{%
      \rIE{1,\FormulaRefAuto{B21TaggedTwoLayerEquality}[thm]}}
    \proofstepwidestar[]{0\in\NatSeg2}{%
      \FormulaRefAuto{PeanoNatSegTwoMembership}[thm]{\rOIa{\rII}}}
    \proofstepwidestar[]{1\in\NatSeg2}{%
      \FormulaRefAuto{PeanoNatSegTwoMembership}[thm]{\rOIb{\rII}}}
    \proofstepwidestar[]{%
      \omega(0,0)=c_0\land\omega(0,1)=c_1}{\FormulaRefAuto{MogiljanskajaDPrimeParametrizationDef}[def]{\rAI{9,10}}
}
    \proofstepwidestar[]{%
      c_0\neq c_1}{%
      \FormulaRefAuto{B21ReserveMarkerDistinct}[thm]}
    \proofstepwidestar[]{%
      \omega\restriction_{I_0}\colon I_0\bij P}{%
      \FormulaRefAuto{RestrictionDef}[def]{6,3},\allowbreak\ \FormulaRefAuto{MogiljanskajaDPrimeDef}[def],\allowbreak\ 8,\allowbreak\ 11,\allowbreak\ 12,\allowbreak\ \FormulaRefAuto{Bijektive Funktion}[def]}
    \proofstepwidestar[]{%
      \nu\colon\mathbb N\times\mathbb N\bij V}{%
      \FormulaRefAuto{MogiljanskajaVParametrizationBijective}[thm]}
    \proofstepwidestar[]{%
      \forall x\in I_1\,\exists!p\in\mathbb N\times\mathbb N\,(x=(1,p))}{%
      \FormulaRefAuto{TaggedProductLayerUniqueRepresentation}[thm],\ 2}
    \proofstepwidestar[]{%
      \forall p\in\mathbb N\times\mathbb N\,(\omega(1,p)=\nu(p))}{%
      \FormulaRefAuto{MogiljanskajaDPrimeParametrizationDef}[def],\ \FormulaRefAuto{x\in A\times B\eqvdash\exists a\in A\,\exists b\in B\,x=(a,b)}[thm]}
    \proofstepwidestar[]{%
      \omega\restriction_{I_1}\colon I_1\bij V}{%
      \FormulaRefAuto{RestrictionDef}[def]{6,3},\ 14,\ 15,\ 16,\ \FormulaRefAuto{Bijektive Funktion}[def]}
    \proofstepwidestar[]{%
      P\cap V=\varnothing}{%
      \FormulaRefAuto{MogiljanskajaReserveCoreFacts}[thm],\allowbreak\ \FormulaRefAuto{MogiljanskajaDPrimeDef}[def],\allowbreak\ \FormulaRefAuto{A\cap(B\cup C)=(A\cap B)\cup(A\cap C)}[thm]}
    \proofstepwidestar[]{%
      \omega\colon I_{D'}\bij D'}{%
      \FormulaRefAuto{CaseFunBijectiveDisjointTargets}[thm]{13,17,5,18},\allowbreak\ \FormulaRefAuto{MogiljanskajaDPrimeDef}[def],\allowbreak\ 3,\allowbreak\ 11,\allowbreak\ 16}
  \closeproofpartwideindR

  \proofpartwide{Bildmenge}
    \proofstepwidestar[]{%
      \omega\colon I_{D'}\bij D'}{%
      \FormulaRefAuto{B21DPrimeFamilyBijective}[thm]}
    \proofstepwidestar[]{%
      \omega\colon I_{D'}\sur D'}{%
      \FormulaRefAuto{Bijektive Funktion}[def]{1}}
    \proofstepwidestar[]{%
      \omega[I_{D'}]=D'}{%
      \FormulaRefAuto{SurjectiveImageEqualsCodomain}[thm]{2}}
  \closeproofpartwide
\end{tabproofsplitwide}

\MogStatementTarget{MogiljanskajaFamilyImages}
\FormulaThmDeltaKR[Mogiljanskaja-Bildmengen]{%
  \begin{aligned}[t]
    &\text{(i)}\quad a[\mathbb N]=L,\\[-2pt]
    &\text{(ii)}\quad b[\mathbb N\times\mathbb N]=D,\\[-2pt]
    &\text{(iii)}\quad \upsilon[\mathbb N\times\mathbb N]=U,\\[-2pt]
    &\text{(iv)}\quad \nu[\mathbb N\times\mathbb N]=V,\\[-2pt]
    &\text{(v)}\quad \omega[I_{D'}]=D'.
  \end{aligned}%
}{%
  \begin{aligned}[t]
    a[\mathbb N]&=L,&
    b[\mathbb N\times\mathbb N]&=D,\\[-2pt]
    \upsilon[\mathbb N\times\mathbb N]&=U,&
    \nu[\mathbb N\times\mathbb N]&=V,\\[-2pt]
    \omega[I_{D'}]&=D'
  \end{aligned}%
}{MogiljanskajaFamilyImages}{%
  \DeltaRow{Mengen}{L\dsep D\dsep U\dsep V\dsep D'\dsep I_{D'}}
  \DeltaRow{Funktionen}{a\dsep b\dsep\upsilon\dsep\nu\dsep\omega}
}
\begin{tabproofsplitwide}
  \proofpartwide{Zusammenstellung der fünf Bildmengen}
    \proofstepwidestar[]{%
      a[\mathbb N]=L}{%
      \FormulaRefAuto{MogiljanskajaLayerLImage}[thm]}
    \proofstepwidestar[]{%
      b[\mathbb N\times\mathbb N]=D}{%
      \FormulaRefAuto{MogiljanskajaLayerDImage}[thm]}
    \proofstepwidestar[]{%
      \upsilon\colon\mathbb N\times\mathbb N\bij U}{%
      \FormulaRefAuto{B21UpsilonBijective}[thm]}
    \proofstepwidestar[]{\upsilon \colon { \mathbb N \times \mathbb N } \to U}{%
      \FormulaRefAuto{Bijektive Funktion}[def]{3}}
    \proofstepwidestar[]{%
      \upsilon[\mathbb N\times\mathbb N]=U}{%
      \FormulaRefAuto{SurjectiveImageEqualsCodomain}[thm]{4}}
    \proofstepwidestar[]{%
      \nu[\mathbb N\times\mathbb N]=V}{%
      \FormulaRefAuto{MogiljanskajaVParametrizationBijective}[thm]}
    \proofstepwidestar[]{%
      \omega[I_{D'}]=D'}{%
      \FormulaRefAuto{MogiljanskajaDPrimeParametrizationBijective}[thm]}
  \closeproofpartwide
\end{tabproofsplitwide}

